\section{Thesis and Dessertation}

        \begin{twocolentry}{
            Jan – Apr 2024
        }
            \href{https://drive.google.com/file/d/1XAY01jkwndByMHBNCZhTZ_OLq4SwbaKe/view?usp=sharing}{\textbf{Optimization Of The Feature Tracking Code To Analyze Solar Quiet Sun Data }} (Masters Thesis)   \href{https://drive.google.com/file/d/177hjcXcHCzlMHXnNjdPv3_OTtgu7IE0A/view?usp=sharing}{\faExternalLink}
            \end{twocolentry}

        \vspace{0.10 cm}
        \begin{onecolentry}
        Guide : \textit{Dr. Anusha L. S., Indian Institute of Astrophysics (IIA), India}\\
            This code identifies, counts and analyzes small scale magnetic features present on the solar photosphere in quiet Sun region (details given in research experience section).
            %Description : Converted an IDL-based feature tracking code into Python and Julia for cross-platform compatibility, optimized its runtime, and developed efficient pipelines to analyze small-scale magnetic features in solar photosphere data from the Sunrise Observatory
        \end{onecolentry}


        \vspace{0.2 cm}

        